%%%%%%%%%%%%%%%%%%%%%%%%%%%%%%%%%
%
%       EXERCÍCIO
%
%%%%%%%%%%%%%%%%%%%%%%%%%%%%%%%%%

\definevalorquestao

\begin{exercicioBanco}[\valorquestao]
Em um certo restaurante, paga-se pelo almoço uma quantia fixa dependendo da escolha feita de prato e bebida. A carne de peixe tem 10\% de preferência, enquanto frango tem 40\% e carne bovina 50\%. As três escolhas bebidas estão condicionadas à opção do prato, segundo a tabela abaixo:
%
\[
    \begin{array}{c|c|c|c}
         & \text{Cerveja} & \text{Água} & \text{Vinho} \\
         \hline
         P(\text{Bebida}\mid\text{Peixe}) & 0,4 & 0,3 & 0,3 \\
         P(\text{Bebida}\mid\text{Frango}) & 0,3 & 0,5 & 0,2 \\
         P(\text{Bebida}\mid\text{Bovina}) & 0,6 & 0,3 & 0,1 \\
    \end{array}
\]
%
Admita os seguintes preços
%
\[
    \begin{array}{c|c|c|c|c|c|c}
        \text{Pedido} & \text{Peixe} & \text{Frango} & \text{Bovina} & \text{Cerveja} & \text{Água} & \text{Vinho} \\ 
        \hline
        \text{Preço} & 12 & 15 & 18 & 6 & 3 & 9
    \end{array}
\]
%
Determine a função de probabilidade para a variável \(X:\) preço do almoço.
\end{exercicioBanco}

